% !TeX root = ../../surface-lagrangian-torus-fibration.tex
% 2026-11-22 / v2.4: 精密主定理，解析入力，証明経路とcutの量化順序．
\section{負 KE 極限を実現する大域 fibration}
\label{ke:global-realization-section}

本節では，負 Kähler--Einstein 計量の Gromov--Hausdorff 極限を，
actual Lagrangian fibration の射影によって実現する主定理を述べる．
後続の節で全域 fibration を構成し，その同じ射影に対する全曲面の距離の上下界を証明する．
最後の統合と対角選択は定理\ref{ke:global-metric-realization}に置く．

\subsection{固定した退化族と計量の規格化}

整数 $d\ge5$ を固定し，まず\eqref{deg:family}の全平面退化を考える．
具体的には，一般位置の平面 $\{L_i=0\}$ と次数 $d$ の式 $F_d$ による
\[
 X_t=\left\{\prod_{i=1}^dL_i+tF_d=0\right\}\subset\PP^3
\]
を扱う．ここで，$F_d$ が三重点で非零，二重線上での零点が
相異なる単純零点であり，十分小さい $t\ne0$ に対して $X_t$ が滑らかであることを仮定する．
族とその幾何学的補助データは固定し，以下で定数の次数一様性や族一様性は要求しない．

$X_t$ 上の負 KE form とその実計量を $\omega_t,g_t$ と書き，
\[
 \operatorname{Ric}(\omega_t)=-\omega_t,\qquad
 [\omega_t]=2\pi(d-4)H,\qquad
 L_t=-\log|t|^2,\qquad R_t=\log L_t
\]
と規格化する．全曲面上の Riemannian distance を $d_{g_t}$ とし，
\[
 \bar g_t=R_t^{-2}g_t,\qquad \bar d_t=R_t^{-1}d_{g_t}
\]
を用いる．$\bar d_t$ は全 $X_t$ 内の経路を許した距離であり，
三重点の内在距離や保護領域内の距離ではない．
極限複体は系\ref{topo:KE-cubical-comparison}の
\[
 \mathcal C_d=
 \{w\in[0,1]^d:\max_i w_i=1,\ |\operatorname{supp}(w)|\le3\}
\]
であり，unit Euclidean squares の貼合せによる全体の長さ距離 $d_\square$ を入れる．
この規格化は，関連 KE 原稿\cite{KELocalCurrent,KECubicalCurrent}および
定理\ref{ke:closed-corner-gh}の規格化と同じである．

\subsection{大域的な幾何学的実現の主定理}

\begin{theorem}[負 KE 極限の Lagrangian fibration による実現]
\label{ke:global-realization}
上の各固定次数と退化族について，実際の marked pants 分解で実現できる
一つの木系 $T$ と，すべての十分小さい $t\ne0$ に対する
\[
 \pi_t:(X_t,\omega_t)\longrightarrow B_d(T)
\]
であって，次の二条件を満たすものが存在する．
\begin{enumerate}
\item $\pi_t$ は定義\ref{def:fibration}の大域 fibration である．
すなわち全 $X_t$ を覆う連結 fiber の連続全射で，正則 fiber は
\emph{actual form $\omega_t$ に関する} Lagrangian $T^2$ である．
特異 fiber には定理\ref{ke:actual-global-fibration}の nodal，branch，core 模型と等方的 strata を用い，
接着の際に whole fibers を失わない．
\item 系\ref{topo:KE-cubical-comparison}の labelled map
$p_T:B_d(T)\to\mathcal C_d$ に対し，$f_t=p_T\circ\pi_t$ と置くと，
\begin{equation}\label{ke:eq:global-realization-distortion}
 \varepsilon_t:=
 \sup_{x,y\in X_t}
 \left|R_t^{-1}d_{g_t}(x,y)
              -d_\square(f_t(x),f_t(y))\right|
 \longrightarrow0
 \qquad(t\longrightarrow0)
\end{equation}
が成り立つ．
\end{enumerate}
木系の組合せ型と component labels は固定するが，nodal values，
collar の位置，作用座標の表示，および基底の parameter 化には $t$ 依存を許す．
任意に既定の $T$ や任意に固定した critical values に対する実現は，
この定理の結論に含めない．
連続性は各固定 $t$ の写像 $\pi_t:X_t\to B_d(T)$ についての条件であり，
$t$ 方向に連続又は滑らかな fibration 族を与えるという主張ではない．
\end{theorem}

第一条件は定理\ref{ke:actual-global-fibration}，第二条件は
定理\ref{ke:global-ambient-lower}と定理\ref{ke:global-ambient-upper}で証明する．
命題\ref{ke:global-base-parameterization}が両評価を同じ $f_t$ に揃え，
定理\ref{ke:global-metric-realization}が共通の遅い対角選択を与える．
次数範囲は各固定 $d\ge5$ であり，正則部分の action--angle structure を超える
全基底の matroidal atlas や special Lagrangian 性は結論に含めない．

\begin{remark}[Atlas との同時実現に十分な値比較]\label{ke:atlas-ke-value-criterion}
$d=5,6$ とし，本定理の $f_t=p_T\pi_t$ に加えて，一つの固定木系 $\widehat T$ に対し，
atlas を持つ別の actual Lagrangian fibration
$\widehat\pi_t:X_t\to B_d(\widehat T)$ を選べたとする．
対応する labelled 収縮を用いて $\widehat f_t=p_{\widehat T}\widehat\pi_t$ と置く．
もし
\[
 \delta_t:=\sup_{x\in X_t}d_\square(f_t(x),\widehat f_t(x))\longrightarrow0
\]
を証明できれば，命題\ref{topo:ke-map-compatibility}の直接の適用により
\[
 \operatorname{dis}(\widehat f_t)
 \le\operatorname{dis}(f_t)+2\delta_t\longrightarrow0
\]
となる．これは新たな比較原理ではなく，既存の一様基底値比較をこの二つの射影へ適用したものである．
命題\ref{ke:moment-relative-replacement}は，この値比較の具体的な十分条件を与える．
同じ保持 corner の全葉，各 component の core 領域，同じ canonical normal
moment と anchor cuts を保ち，marked product/cone 商と全 overlap の
whole-fiber 接着を持つ相対変更なら，
$\delta_t\le C(1+\log(N_{j(t)}+1))/R_t$ となる．
Reduced 方向の変更を微分まで小さくする必要はない．
定理\ref{ng8:componentwise-critical-area}と系\ref{ng8:p0-small-area-flexibility}により，
相対面積実現の有限判定と source を保持する微小変更は得られた．
Actual core chart には命題\ref{ke:actual-anchor-fan-criterion}の
anchor period 条件を使える．五次の固定 anchor では，
系\ref{ke:k4-protected-anchor-selection}が保持 tail の period による有限判定を与える．
未解決なのは，actual データにおけるその条件の検証，所定の atlas targets に対する
面積配分の可解性と，反対側 neck の切片・整数 framing を含む全 open overlaps の一致である．
定理\ref{main:action-atlas}の各固定 $t$ における存在だけではこの条件は得られず，
一般の symplectic conjugacy だけで critical-action の障害が消えるとも主張しない．
\end{remark}

$f_t$ は全射なので，\eqref{ke:eq:global-realization-distortion}により，
その graph correspondence，または命題\ref{topo:ke-map-compatibility}により
\begin{align*}
 d_{\mathrm{GH}}\bigl((X_t,\bar d_t),(\mathcal C_d,d_\square)\bigr)
 &\le\tfrac12\varepsilon_t,\\
 \sup_{c\in\mathcal C_d}
 \operatorname{diam}_{\bar d_t}f_t^{-1}(c)
 &\le\varepsilon_t,\\
 \sup_{b\in B_d(T)}
 \operatorname{diam}_{\bar d_t}\pi_t^{-1}(b)
 &\le\varepsilon_t
\end{align*}
を同時に得る．ここで直径は全空間の距離を制限した ambient diameter である．
特異 fiber もこの上限に含み，元の尺度での収縮や fiber 内在直径の収縮までは要求しない．
$d=5$ では $p_T$ は同相である．$d\ge6$ では，
一つの $p_T$-fiber 上にある $\pi_t$-fibers の集合全体が縮むことが必要であり，
各 torus の点ごとの収縮だけでは十分でない．

\begin{remark}[径を一にする規格化との関係]
\label{ke:global-realization-scale}
$D_d=\operatorname{diam}(\mathcal C_d,d_\square)>0$ とする．
\eqref{ke:eq:global-realization-distortion}と全射性から
\[
 \left|R_t^{-1}\operatorname{diam}(X_t,g_t)-D_d\right|
 \le\varepsilon_t
\]
である．従って径を一に規格化した計量の極限は
$(\mathcal C_d,D_d^{-1}d_\square)$ となる．
定理\ref{topo:diameter-six-labels}により $d\ge6$ では $D_d=2\sqrt5$ であり，
五次では既存の有限計量計算\cite[\texttt{thm:intro-diameter-C5}]{KECubicalCurrent}から
$D_5=\sqrt{13}$ である．これらの有限計量結果と主定理から，
$R_t^{-1}\operatorname{diam}(X_t,g_t)$ の極限も定まる
（系\ref{ke:global-diameter-limit}）．
元の負 KE 尺度の極限，$R_t$ による規格化，および径一の規格化は区別する．
\end{remark}

\subsection{解析入力の出典と適用範囲}
\label{ke:global-realization-transfer}
\label{ke:proof-inputs}

主定理の証明では，退化族の局所解析，固定成分の cusp 幾何，および本稿の
actual fibration の構成を次のように接続する．退化族に関する局所入力の正本は
\cite[Part~IV]{KELocalCurrent}であり，以下に英字で記す定理ラベルは同稿を指す．
各固定 $d$ と族に対する定数を使い，固定 compact 上の収束から
任意の moving 領域における一様性を推論することはない．

\begin{enumerate}
\item \emph{Global volume 下界と oriented 領域の volume 上界．}
\texttt{coarse:cor:degree-d-volume-lower} は，全曲面上の滑らかな capped density に
対する actual KE 体積の下界を与える．
\texttt{coarse:lem:pair-arm-density} と
\texttt{coarse:lem:complete-volume-upper} は，source を含む oriented arm で
比較する model density と上界を与える．後者の完全計量比較で必要な
log volume ratio の上方有界性は，各固定 $t$ の閉包と model の人工境界での発散から
検証する条件であり，未構成の actual metric buffer を仮定しない．
これらを用い，本稿で別に証明する
補題\ref{ke:capped-projective-lower}の固定残余領域上の射影計量下界，
補題\ref{ke:source-tube-potential}の正の potential と
補題\ref{ke:positive-poisson-log-gradient}の局所対数勾配，
命題\ref{ke:source-arm-uniform-tensor}の三境界 Schwarz 評価を組み合わせる．
体積密度の二側有界性だけを tensor 二側比較に置き換えてはいない．

\item \emph{全深部 triple の tensor 比較と complete simplex の幾何．}
\texttt{coarse:thm:full-triple-tensor} は，三重点の全 phase と全深度比を含む
$\min_a\ell_a\ge m$ 上で，$L_t\ge L_*$，$m_*\le m<L_t/3$ を満たす全 cut に一様な
$O(m^{-1})$ の tensor 比較を全固定微分階で与える．
その局所的な核は \texttt{coarse:thm:boundary-to-triple} である．
同定理は固定 face の volume 下界と，各固定 $t$ の radial 閉包近傍での
smooth 負 KE 計量という入力を分離し，全平面退化への適用では
上の capped 下界でこれを満たす．
Volume ratio が $1$ に近づく評価と
Savin\cite[Theorem~1.3]{SavinSmall}の小摂動正則性を組み合わせる段階が，
有界な volume ratio のみからの推論との相違である．
\texttt{simplex:lem:rank-one-schwarz}，
\texttt{simplex:prop:dominant-product}，
\texttt{simplex:cor:active-cogram}，
\texttt{simplex:thm:closed-corner-model} は，同じ complete simplex 上の
長さと active covectors の比較を与える．
本稿の定理\ref{ke:slow-cut-closed-corner-gh}はこれらを
actual Lagrangian な closed corner と有限 cut の距離評価へ移す．
全曲面の source 又は二重曲線の全領域が triple 比較の適用域に含まれるとはしない．

\item \emph{固定 residual compact と固定 end の normal moment．}
\texttt{pair:prop:bundle-conic}，
\texttt{pair:cor:global-primitive}，
\texttt{pair:thm:fixed-end-moment}，
\texttt{pair:cor:actual-common-moment} は，source-free な固定残余基底の
正則 bundle 座標，canonical primitive，固定 end における moment の比較を与える．
固定 compact と end を先に取り，その後 $t\to0$ とする順序を用いる．
Moving な全軌道を含む actual flow buffer，source の全 moment 接続，
複数 end の同時接着は，本稿の
節\ref{ke:staged-neck-section}，
\ref{ke:neck-corner-exhaustion-section}および
定理\ref{ke:source-all-moment-exit}で追加して証明する．
\texttt{pair:lem:diagonal-fixed-end} の遅い対角選択は，
これらの評価を証明した後に使う原理である．

\item \emph{固定成分上の smooth convergence と SNC cusp 評価．}
Song\cite[Theorem~1.2]{SongNegativeKE}による成分極限の同定と
fixed-gauge convergence は，上記
\texttt{pair:thm:fixed-end-moment} の証明で対象族へ適用される．
成分は $\PP^2$ から $d-1$ 本の直線を除いたもので，
log canonical bundle は $\mathcal O_{\PP^2}(d-4)$ である．
Rochon--Zhang\cite[\S5, (5.3)--(5.6), Theorem~5.1]{RochonZhang}の
bounded cusp derivatives を持つ potential を用い，
本稿の節\ref{ke:component-core-geometry-section}と
\ref{ke:component-core-transfer-section}で正の residue，
固定 compact の正則化と小さい actual 移送を得る．
Divisor が smooth な場合だけの polyhomogeneous expansion は使わず，
$t$ とともに動く任意の compact 上での一様収束も仮定しない．
\end{enumerate}

一般の Moser 同定が Lagrangian 性を運ぶことと，その移送が KE 距離を保つことは別である．
本稿では下界と上界を actual $g_t$ に対して直接証明する．
\cite{KECubicalCurrent}の
$\mathrm{(G)}_d+\mathrm{(CP)}_d+\mathrm{(PBE)}_d$ による条件付き GH 定理，
又は $\mathrm{(TC)}_5$ を主定理の入力にはしない．
既存の近似写像 $\Phi_t$ との一様な値比較から
$f_t=p_T\pi_t$ へ距離歪みを移す命題\ref{topo:ke-map-compatibility}は，
別の有効な十分条件として保持する．
本稿の直接証明を，これらの外部比較仮説自体の検証へ読み替えることはできない．

\subsection{同じ射影を保つ六段階の証明}
\label{ke:proof-route}

構成を確認した後に，同じ射影の距離評価を追うための経路を以下にまとめる．
局所 GH 評価は第3段階，全曲面の ambient 距離比較は第5段階である．

\begin{enumerate}
\item \emph{Neck と closed corner の全葉を一致させる．}
系\ref{ke:neck-corner-saturated-union}は，
共通の向き付き normal action と canonical moment を保ち，
全ての実際の overlap で reduced leaves と whole fibers を一致させる．
補題\ref{ke:fixed-marked-tree-exhaustion}により，
実際の marked pants 分解から得る木系を全段階と全ての小さい複素 $t$ で固定する．

\item \emph{Actual core を埋め，全曲面の fibration を作る．}
定理\ref{ke:actual-global-fibration}は，closed moment links，
primitive の相対不変化と nontrapping，saddle を動かす葉族補間，
補間後の全 fiber 単位の trim を用いる．
有限な closed whole-fiber quotient の Hausdorff 性を先に確保し，
全 $X_t$ の被覆，連結 fibers，actual Lagrangian 性と特異模型を得る．
Source の pinched tori と保持 corner の葉はこの相対構成で保たれる．

\item \emph{後続の接続に使う局所費用を評価する．}
定理\ref{ke:source-all-moment-exit}は，source の全共通 moment で，
同じ canonical moment の残余境界まで actual 長さ $O(1)$ の退出を与える．
Source 軸を含む arm tensor 比較と中央点からの短い直接退出による証明なので，
Gaussian 可積分性を必要としない．
定理\ref{ke:slow-cut-closed-corner-gh}で closed corner の距離を評価し，
命題\ref{ke:component-logarithmic-connectors}で各 component の
core・変更 collar・neck 尾部の費用を $C(1+\log(N_j+1))$ に抑える．

\item \emph{一つの基底 parameter 化を選ぶ．}
命題\ref{ke:global-base-parameterization}により，第2段階の同じ fiber partition を
$B_d(T)$ と同定し，$f_{j,t}=p_T\pi_{j,t}$ とする．
保持 corner の像，neck strips，component sectors をこの写像で比較し，
内部木の全枝にも一つの canonical moment を用いる．

\item \emph{この同じ写像の全域下界と上界を証明する．}
下界は定理\ref{ke:sharp-global-one-active}の source を除かない scalar 障壁から，
系\ref{ke:global-depth-lipschitz}の global clipped-depth map，
命題\ref{ke:global-depth-base-comparison}の一様な基底値比較を経て，
定理\ref{ke:global-ambient-lower}へ至る．
長さ評価は全曲面の任意の経路に成立し，領域の出入りの回数に依存しない．
上界は補題\ref{ke:coherent-corner-portals}の共通 moment の接続を
高々 $3\binom{\min\{d,6\}}3\le60$ 本の square gallery に沿って結ぶ
定理\ref{ke:global-ambient-upper}で得る．
下界の値比較で用いる接続は第3段階で証明してあり，
全域距離上界を先取りしない．

\item \emph{有限段階の誤差を一つの対角選択にまとめる．}
定理\ref{ke:global-metric-realization}は，
式\eqref{ke:eq:global-realization-explicit-cut-budget}の両側評価に
共通の遅い選択 $j=j(t)$ を適用し，
\eqref{ke:eq:global-realization-distortion}を結論する．
各 torus だけでなく，$p_T$ が潰す木上の全 fiber 合併を同じ評価に含める．
\end{enumerate}

\subsection{切断尺度と退化 parameter の量化順序}
\label{ke:cut-selection-order}

固定族の下で，各有限段階 $j$ のデータを次の順に選ぶ．
これは任意に先指定した moving 領域への一様評価ではなく，
既に証明した局所評価を同じ全域構成で用いるための選択である．
\begin{enumerate}
\item $b_j\to\infty$ と $\varepsilon_j\downarrow0$ を取り，
旧 corner cut $m_j=6b_j$，固定 residual compact，既存の normal cut $n_j$，
共通 moment 幅 $\iota_j$ を順に選ぶ．
固定 residual domain の接続定数 $C_j^{\rm neck}$ は，次の $N_j$ より先に固定する．
\item Component cut $N_j$ を，既存の下限，
$N_j\iota_j\gg1$，および
\[
 C_j^{\rm neck}\le1+\log(N_j+1)
\]
を満たすように選ぶ．これにより core と変更 collar を作り，
neck と component の全接続に一つの対数的予算を使える．
\item Scalar 障壁専用の補助 cut $\mathfrak n_j$ と境界定数を先に決め，
その後に lower clipping depth $\mathfrak m_j$ と upper factor $K_j$ を選ぶ．
$\mathfrak m_j$ は固定段階の core の深度上界も支配させる．
$\mathfrak n_j$ と既存の $n_j$，$\mathfrak m_j$ と旧 corner cut $m_j$ は別の量である．
正確な下限は定理\ref{ke:sharp-global-one-active}と
命題\ref{ke:global-depth-base-comparison}に置く．
\item 最後に，全ての構成条件と誤差評価に加え，
\[
 R_t\ge j\{1+\log\mathfrak m_j+\log K_j+\log(N_j+1)\}
\]
を保証する閾値 $r_j>0$ を選ぶ．
$r_{j+1}<r_j/2$ として，$j(t)=\max\{j:|t|<r_j\}$ と置く．
\end{enumerate}
残余定数の増大率や固定 compact の収束閾値は有効に評価していないので，
この量化順序は普遍的な収束速度を与えない．また，各 $\pi_t$ の連続性と
$t$ 方向の連続又は滑らかな選択は区別する．

Source の marking には命題\ref{ke:source-marking}の条件が必要である．
指定 nodal fiber が actual source core と交わり，候補近似 $\Phi_t$ が raw map と
一様に $o(1)$ で近いとする．このとき
$\sup_{X_t}d_\square(\Phi_t,p_T\pi_t)\to0$ を得るには，
その nodal value $b_{ij,t}$ が $p_T(b_{ij,t})\to e_i+e_j$ を満たす必要がある．
系\ref{ke:neck-base-value-control}の構成はこの位置条件を満たす．
任意の先指定 critical values 又は任意の第二 periods の実現を，
以上の六段階に追加して主張するものではない．
